\chapter{Requirement Specification}

\section{Existing System}

Legacy compliance management systems rely on manual spreadsheet auditing and fragmented Software-as-a-Service (SaaS) platforms. Governance teams attempt to map regulatory standards, such as ISO/IEC 27001:2022, to technical controls using decentralized, static documents. This approach introduces human error, lacks version control, and fails to maintain referential integrity across overlapping frameworks. The manual tracking of configuration states results in heavy operational labor expenditures and administrative burden.

Existing point-in-time auditing solutions perform delayed security assessments, meaning data collected from internal network systems becomes stale upon generation. Legacy platforms employ generic software flows lacking specialized relational constraints, leading to latency when processing endpoint telemetry and an inability to maintain determinism in vulnerability resolution. The reliance on proprietary SaaS architectures introduces data sovereignty risks, because internal compliance metrics and infrastructure configurations are transmitted to multi-tenant public cloud environments.

\section{Proposed System}

To overcome the limitations of legacy spreadsheet mapping and asynchronous SIEM aggregation, the ICMS platform introduces a centralized, deterministic orchestration engine. The system leverages a multi-tier architectural topology designed to maintain separation of concerns while providing compliance tracking within a single-tenant, isolated environment.

\subsection{Technology Stack Selection}

Selecting the Python-based Django framework for the backend application layer provides an environment characterized by an integrated Object-Relational Mapper (ORM) and native security middleware. The Django ORM facilitates the implementation of performance optimizations, specifically $O(1)$ set-based database aggregations. By utilizing query projections and conditional filtering at the ORM level, the system ensures high-volume telemetry remains computationally manageable. The Django REST Framework (DRF) extends these capabilities by providing modular serialization pipelines and support for authentication schemas, including stateless JSON Web Token (JWT) rotation.

The selection of React 18 and the Vite build tool establishes a Single Page Application (SPA) architecture designed for data visualization. React's Virtual DOM provides an advantage in the context of security monitoring, allowing the interface to render compliance state changes without causing browser thread execution delays. This reactive rendering ensures analysts maintain an operational experience when viewing dense matrices of SIEM alerts. Vite enhances the development and production lifecycle by utilizing native ES modules and Hot Module Replacement (HMR), reducing build latency compared to legacy Webpack-based bundling strategies.

PostgreSQL serves as the primary data persistence tier, providing the ACID compliance required for maintaining the integrity of organizational compliance records. The relational engine handles the concurrent read/write loads generated by SIEM telemetry ingestion and automated framework mapping. By utilizing B-Tree indexing and relational constraints, PostgreSQL guarantees cross-framework compliance metrics remain accurate. This data integrity satisfies the evidence requirements of ISO 27001 and SOC 2 audits, where the historical state of security controls must remain verifiable and immutable.

To ensure the platform remains responsive during data-heavy operations, the architecture incorporates Celery and Redis as the asynchronous orchestration layer. Decoupling I/O-bound tasks, such as the periodic synchronization of telemetry from the Wazuh SIEM REST API, from the main Django application thread prevents HTTP request timeouts and maintains dashboard availability. Redis acts as a message broker, facilitating the transport of task payloads between the web server and the background worker fleet. This distributed task execution model allows the platform to scale its ingestion capacity without impeding the interaction capabilities of the primary API service.

\subsection{High-Level Architecture}

The Presentation Tier consists of a React-based SPA that manages the user session and provides an interface for compliance monitoring. This tier communicates with the Application Logic Tier via stateless JSON payloads over an encrypted HTTPS channel. By offloading the state management and rendering logic to the client-side browser environment, the system ensures central server resources remain available for data processing and framework mapping tasks.

The Application Logic Tier, powered by the Django REST Framework, serves as the gateway for compliance orchestration. This tier manages the Identity and Access Management (IAM) lifecycle, including the generation and rotation of JSON Web Tokens (JWT). Interactions with the Data Persistence Tier occur via TCP connections managed by the Django ORM, which enforces relational integrity and executes query aggregations. Asynchronous task execution relies on a distributed architecture incorporating Celery and a Redis-based message broker. The system utilizes a Pub/Sub messaging model to enqueue compliance synchronization tasks, allowing Celery workers to process telemetry data from the external Wazuh SIEM Manager out-of-band. Figure \ref{fig:detailed_architecture} maps the protocol flow and multi-tier layout.

\begin{figure}[H]
    \centering
    \resizebox{\textwidth}{!}{
    \begin{tikzpicture}[
        service/.style={rectangle, draw, minimum width=5cm, minimum height=2cm, align=center, fill=blue!5, thick},
        database/.style={cylinder, draw, shape border rotate=90, minimum width=3cm, minimum height=2.5cm, align=center, fill=gray!10, thick},
        external/.style={rectangle, draw, dashed, minimum width=5cm, minimum height=2cm, align=center, fill=red!5, thick},
        arrow/.style={-Stealth, thick},
        label_node/.style={fill=white, inner sep=3pt, font=\small\itshape}
    ]
        % Nodes with Absolute Coordinates
        \node[service] (frontend) at (0,0) {\textbf{React SPA (Frontend)}\\[2mm] \small $\bullet$ Port 80/443 (HTTP/S)\\ $\bullet$ Vite Bundler\\ $\bullet$ Axios Interceptors};
        
        \node[service] (backend) at (8,0) {\textbf{Django Core API}\\[2mm] \small $\bullet$ Port 8000 (WSGI)\\ $\bullet$ DRF Routers\\ $\bullet$ JWT Middleware\\ $\bullet$ $O(1)$ Query Engine};
        
        \node[database] (db) at (15,3) {\textbf{PostgreSQL (Data Tier)}\\[2mm] \small $\bullet$ Port 5432 (TCP)\\ $\bullet$ Relational DB\\ $\bullet$ ACID Compliance};
        
        \node[service] (redis) at (15,-2) {\textbf{Redis (Message Broker)}\\[2mm] \small $\bullet$ Port 6379 (TCP)\\ $\bullet$ In-Memory Cache\\ $\bullet$ Task Queuing};
        
        \node[service] (celery) at (15,-7) {\textbf{Celery (Async Tier)}\\[2mm] \small $\bullet$ Worker Nodes\\ $\bullet$ Task Scheduling\\ $\bullet$ Background Sync};
        
        \node[external] (wazuh) at (8,-8) {\textbf{Wazuh Manager (External)}\\[2mm] \small $\bullet$ Port 55000\\ $\bullet$ SCA REST API};

        % Environment Bounding Boxes
        \begin{scope}[on background layer]
            \node[draw, blue!60, dashed, thick, fill=blue!2, inner xsep=20pt, inner ysep=35pt, fit=(frontend)] (client_env) {};
            \node[font=\bfseries\color{blue!80}, anchor=south, yshift=5pt] at (client_env.north) {Client Environment};
            
            \node[draw, gray!80, thick, fill=gray!10, inner xsep=20pt, inner ysep=35pt, fit=(backend) (db) (redis) (celery)] (server_env) {};
            \node[font=\bfseries\color{gray!80}, anchor=south, yshift=5pt] at (server_env.north) {Isolated Docker Bridge Network};
        \end{scope}

        % Edge Routing
        \draw[arrow] (frontend.east) -- (backend.west) 
            node[midway, above, yshift=2pt, font=\small\itshape] {HTTPS (JSON)};
            
        \draw[arrow] ([yshift=8mm]backend.east) |- (db.west) 
            node[pos=0.75, above, font=\small\itshape] {ORM (TCP)};
            
        \draw[arrow] ([yshift=-8mm]backend.east) |- (redis.west) 
            node[pos=0.75, above, font=\small\itshape] {AMQP};
            
        \draw[arrow] (redis.south) -- (celery.north) 
            node[midway, right, font=\small\itshape] {Task Payload};
            
        \draw[arrow] (celery.east) -- ++(1,0) |- (db.east) 
            node[pos=0.75, above, font=\small\itshape] {Async Write};
            
        \draw[arrow] ([xshift=-1.5cm]backend.south) -- ([xshift=-1.5cm]wazuh.north) 
            node[midway, left, font=\small\itshape] {API Polling};
            
        \draw[arrow] (celery.south) -- (15, -10) -- (8, -10) 
            node[midway, above, font=\small\itshape] {Telemetry Sync} -- (wazuh.south);

    \end{tikzpicture}
    }
    \caption{Refined N-Tier System Architecture and Protocol Flow}
    \label{fig:detailed_architecture}
\end{figure}

\subsection{Dockerized Infrastructure and Network Isolation}

The implementation of a Dockerized virtual bridge network establishes a perimeter that isolates the data persistence tier from external network environments. By restricting the PostgreSQL and Redis services to an internal, non-routable network, the architecture ensures these components remain inaccessible to host-level scans. Managing data persistence across container lifecycles requires Docker volumes to prevent the loss of compliance telemetry. Orchestrating the system startup sequence through deterministic health-check dependencies ensures the platform remains resilient against race conditions during deployment. The configuration defines execution priorities, where the Django API container monitors the execution of system-level commands before initiating database migrations or starting the WSGI server. Figure \ref{fig:docker_mesh} depicts this multi-tier isolation and health monitoring network.

\begin{figure}[H]
    \centering
    \resizebox{\textwidth}{!}{
    \begin{tikzpicture}[
        container/.style={draw, black!80, thick, fill=white, text=black, align=center, minimum width=4.5cm, minimum height=1.8cm, rounded corners},
        host_port/.style={font=\bfseries\small, align=center, text=blue!80},
        health_arrow/.style={-Stealth, dashed, red!70, thick},
        expose_arrow/.style={-Stealth, thick, blue!70}
    ]
        % Nodes
        \node[container] (react) at (0, 0) {\textbf{React SPA Container}\\[1.5mm] \scriptsize Vite Server (Exposed 80/443)};
        \node[container] (django) at (0, -3) {\textbf{Django API Container}\\[1.5mm] \scriptsize Gunicorn (Exposed 8000)};
        \node[container] (celery) at (0, -6) {\textbf{Celery Container}\\[1.5mm] \scriptsize Background Daemon};
        \node[container] (db) at (7, -1.5) {\textbf{PostgreSQL Container}\\[1.5mm] \scriptsize Internal Port 5432};
        \node[container] (redis) at (7, -4.5) {\textbf{Redis Container}\\[1.5mm] \scriptsize Internal Port 6379};
    
        % Host Ports
        \node[host_port] (port80) at (-5, 0) {Host Interface\\(Port 80/443)};
        \node[host_port] (port8000) at (-5, -3) {Host Interface\\(Port 8000)};
    
        % Exposure Routing
        \draw[expose_arrow] (port80) -- (react);
        \draw[expose_arrow] (port8000) -- (django);
    
        % Health-check dependencies
        \draw[health_arrow] (django.east) -- ++(0.5,0) |- (db.west) node[pos=0.7, above, font=\scriptsize\color{red!80}, yshift=2pt] {depends\_on};
        \draw[health_arrow] (django.east) -- ++(0.5,0) |- ([yshift=2mm]redis.west) node[pos=0.7, above, font=\scriptsize\color{red!80}, yshift=2pt] {depends\_on};
        \draw[health_arrow] (celery.east) -- ++(0.5,0) |- ([yshift=-2mm]redis.west) node[pos=0.7, below, font=\scriptsize\color{red!80}, yshift=-2pt] {depends\_on};
    
        % Bounding Box for Docker Network
        \begin{scope}[on background layer]
            \node[draw, black!70, ultra thick, dashed, fill=gray!10, inner xsep=20pt, inner ysep=30pt, fit=(react) (django) (celery) (db) (redis)] (docker_env) {};
            \node[font=\bfseries\color{black!80}, anchor=south, yshift=5pt] at (docker_env.north) {Docker Daemon (Service Mesh)};
        \end{scope}
    
        % Legend
        \begin{scope}[shift={(3.5, -9.0)}]
            \node[draw, rectangle, fill=white, inner sep=10pt, minimum width=8.5cm, minimum height=1.5cm] (legend_box) at (0,0) {};
            \draw[expose_arrow] (-3.7, 0.2) -- (-2.2, 0.2) node[right, text=black, font=\small] {External Port Exposure};
            \draw[health_arrow] (-3.7, -0.2) -- (-2.2, -0.2) node[right, text=black, font=\small] {Health Dependency (service\_healthy)};
        \end{scope}
    \end{tikzpicture}
    }
    \caption{Docker Service Mesh: Port Exposure and Health-Check Dependencies}
    \label{fig:docker_mesh}
\end{figure}

\section{Requirements Specification}

Requirement engineering for the ICMS platform involves mapping regulatory mandates to technical specifications. Aligning the platform with ISO/IEC 27001:2022 and SOC 2 Type II standards requires the translation of governance objectives into deterministic functional capabilities. This process ensures architectural components satisfy compliance or operational constraints. The requirement analysis phase establishes the foundation for a system capable of maintaining visibility across distributed infrastructure by prioritizing automated telemetry ingestion and data aggregation.

\section{Functional Requirements}

The functional requirements define the capabilities the ICMS architecture must possess to execute its core operations:

\begin{itemize}
    \item The system must enforce a 3-tier Role-Based Access Control (RBAC) model integrated with Time-Based One-Time Password (TOTP) secondary authentication.
    \item The platform must automate the ingestion of endpoint security telemetry from the Wazuh SCA module via authenticated REST API synchronization.
    \item The system must resolve conflicting telemetry states using a deterministic Worst-Case Priority state-mutation algorithm to ensure data integrity.
    \item The platform must generate server-side compliance reports in PDF format utilizing CSS2.1 compliant layouts for audit documentation.
    \item The system must provide an administrative interface for the dynamic management of frameworks, controls, and Wazuh SIEM mapping identifiers.
    \item The application must implement a stateless session management lifecycle utilizing JSON Web Token (JWT) rotation for all authenticated requests.
\end{itemize}

\section{Non-Functional Requirements}

The non-functional requirements dictate the operational constraints and performance thresholds required to maintain system reliability:

\begin{itemize}
    \item The system must isolate the PostgreSQL persistence tier and Redis message broker within a private Dockerized virtual bridge network.
    \item The backend must utilize $O(1)$ set-based query aggregations to ensure dashboard rendering latency remains below 200 milliseconds.
    \item The architecture must employ asynchronous background orchestration via Celery and Redis to maintain non-blocking UI responsiveness during data ingestion.
    \item The frontend must utilize a responsive Flexbox and Cyber-Card based architecture to mitigate analyst alert fatigue and reduce cognitive load.
    \item The platform must ensure high availability through the implementation of Celery task retry logic and persistent message queuing for all telemetry sync tasks.
    \item The system must maintain a granular audit logging engine that records all administrative actions and automated system events for forensic traceability.
\end{itemize}

\section{User Requirements}

The operational efficacy of the ICMS platform relies on delivering a user experience tailored to its primary actors: the System Administrator and the Security Auditor. The System Administrator requires a control plane capable of managing infrastructure settings without exposing database commands. They must possess the capability to configure the synchronization lifecycle, manage RBAC privileges across user accounts, and override automated telemetry when necessary. The interface must present these administrative controls with a low cognitive load, utilizing visual feedback through the component library to confirm successful CRUD operations within the Django ORM.

The Security Auditor interacts with the platform as an analytical tool to assess regulatory adherence. Their primary requirement is parsing compliance data without the visual clutter associated with legacy GRC software. The auditor requires an RBAC boundary that prevents modification of infrastructure mapping while offering access to historical scan logs and PDF export functionalities. The React SPA architecture must guarantee the dashboard remains responsive during concurrent data ingestion from the Wazuh SIEM, employing memoization to ensure DOM updates do not introduce browser thread execution delays.

\section{System Models and Components}

This section outlines the systemic data exchange boundaries and flow orchestrations utilized by the platform to synchronize external SIEM telemetry with the internal compliance tracking engine.

\subsection{Context Diagram}

The Context Diagram defines the interaction between the ICMS core logic and external entities, illustrating the boundaries of the application and the data streams required for operation. Figure \ref{fig:context_diagram} maps the interaction flow from human operators to infrastructure components.

\begin{figure}[H]
    \centering
    \resizebox{\textwidth}{!}{
    \begin{tikzpicture}[
        node distance=3cm,
        process/.style={circle, draw=blue!80!black, thick, fill=blue!10, minimum size=3.5cm, align=center, font=\bfseries},
        entity/.style={rectangle, draw=gray!80!black, thick, fill=gray!10, minimum width=3.5cm, minimum height=1.5cm, align=center, font=\bfseries},
        flow/.style={-Stealth, thick, draw=black!80},
        flowlabel/.style={font=\scriptsize, align=center, text=black!90}
    ]
        % Main Process (Center)
        \node[process] (icms) at (0,0) {0.0\\ICMS\\Platform};
        
        % External Entities (Left, Right, Top)
        \node[entity] (wazuh) at (-7,0) {Wazuh SIEM\\Manager};
        \node[entity] (auditor) at (7,0) {Security\\Auditor};
        \node[entity] (admin) at (0, 4.5) {System\\Administrator};
        
        % Data Flows: Wazuh <-> ICMS
        \draw[flow] ([yshift=3mm]wazuh.east) -- ([yshift=3mm]icms.west) node[midway, above, flowlabel] {Raw Security Telemetry};
        \draw[flow] ([yshift=-3mm]icms.west) -- ([yshift=-3mm]wazuh.east) node[midway, below, flowlabel] {SCA Telemetry Request};
        
        % Data Flows: Auditor <-> ICMS
        \draw[flow] ([yshift=-3mm]auditor.west) -- ([yshift=-3mm]icms.east) node[midway, below, flowlabel] {Report Query};
        \draw[flow] ([yshift=3mm]icms.east) -- ([yshift=3mm]auditor.west) node[midway, above, flowlabel] {Compliance Dashboards\\\& PDF Reports};
        
        % Data Flows: Admin <-> ICMS
        \draw[flow] ([xshift=-3mm]admin.south) -- ([xshift=-3mm]icms.north) node[midway, left, flowlabel, align=right] {Configuration Parameters\\\& Sync Commands};
        \draw[flow] ([xshift=3mm]icms.north) -- ([xshift=3mm]admin.south) node[midway, right, flowlabel, align=left] {Audit Logs\\\& System Alerts};
        
        % System Boundary Context
        \begin{scope}[on background layer]
            \node[draw, dashed, black!50, thick, rounded corners, fit=(icms) (wazuh) (auditor) (admin), inner sep=1cm] (boundary) {};
            \node[anchor=north west, font=\itshape\color{black!60}] at (boundary.north west) {Organizational Context Boundary};
        \end{scope}
    \end{tikzpicture}
    }
    \caption{Context Diagram: High-Level System Interactions}
    \label{fig:context_diagram}
\end{figure}

\subsection{Data Flow Diagram}

The Level-1 Data Flow Diagram (DFD) offers a visualization of how data payloads traverse the platform's internal architecture, mapping the transformation of raw SIEM logs into structured compliance metrics. Figure \ref{fig:dfd_level_1} outlines the deterministic pipeline executed by the ICMS.

\begin{figure}[H]
    \centering
    \resizebox{\textwidth}{!}{
    \begin{tikzpicture}[
        node distance=2.5cm,
        process/.style={circle, draw=blue!80!black, thick, fill=blue!10, minimum size=2.8cm, align=center, font=\bfseries\small},
        entity/.style={rectangle, draw=gray!80!black, thick, fill=gray!10, minimum width=3.5cm, minimum height=1.5cm, align=center, font=\bfseries\small},
        datastore/.style={rectangle, draw=gray!80!black, thick, fill=white, minimum width=3.5cm, minimum height=1.2cm, align=center, font=\bfseries\small, path picture={\draw[thick] ([xshift=12pt]path picture bounding box.north west) -- ([xshift=12pt]path picture bounding box.south west);}},
        flow/.style={-Stealth, thick, draw=black!80},
        flowlabel/.style={font=\scriptsize, align=center, text=black!90, inner sep=0pt}
    ]
        % External Entities (Left)
        \node[entity] (admin) at (-6.5, 4.5) {System\\Administrator};
        \node[entity] (wazuh) at (-6.5, 1.5) {Wazuh SIEM\\Manager};
        \node[entity] (auditor) at (-6.5, -4.5) {Security\\Auditor};

        % Processes (Center)
        \node[process] (p1) at (0, 4.5) {1.0\\Manage\\Access};
        \node[process] (p2) at (0, 1.5) {2.0\\Ingest\\Telemetry};
        \node[process] (p3) at (0, -1.5) {3.0\\Map\\Controls};
        \node[process] (p4) at (0, -4.5) {4.0\\Generate\\Reports};

        % Data Stores (Right)
        \node[datastore] (d1) at (6.5, 4.5) {D1: User DB};
        \node[datastore] (d3) at (6.5, 1.5) {D3: Telemetry DB};
        \node[datastore] (d2) at (6.5, -1.5) {D2: Framework DB};

        % Flows: Manage Access
        \draw[flow] ([yshift=4mm]admin.east) -- ([yshift=4mm]p1.west) node[midway, above=3pt, flowlabel] {Auth Credentials};
        \draw[flow] ([yshift=-4mm]p1.west) -- ([yshift=-4mm]admin.east) node[midway, below=3pt, flowlabel] {JWT Session Token};
        \draw[flow] (p1.east) -- (d1.west) node[midway, above=2pt, flowlabel] {Verification Data};

        % Flows: Ingest Telemetry
        \draw[flow] (wazuh.east) -- (p2.west) node[midway, above=2pt, flowlabel] {Raw SCA JSON Data};
        \draw[flow] (p2.east) -- (d3.west) node[midway, above=2pt, flowlabel] {Parsed Endpoint State};

        % Flows: Map Controls
        \draw[flow] (d2.west) -- (p3.east) node[midway, above=2pt, flowlabel] {Framework Rules};
        \draw[flow] (p3.35) -- (d3.190) node[midway, sloped, above=3pt, flowlabel] {Mapped Compliance State};
        \draw[flow] (d3.210) -- (p3.15) node[midway, sloped, below=3pt, flowlabel] {Read Unmapped State};

        % Flows: Generate Reports
        \draw[flow] ([yshift=4mm]auditor.east) -- ([yshift=4mm]p4.west) node[midway, above=3pt, flowlabel] {Report Parameters};
        \draw[flow] ([yshift=-4mm]p4.west) -- ([yshift=-4mm]auditor.east) node[midway, below=3pt, flowlabel] {PDF Audit Report};
        
        % Long-distance route from D3 to P4
        \draw[flow] (d3.east) -- ++(1.5,0) |- node[pos=0.25, right=3pt, flowlabel] {Aggregated Metrics} (p4.east);

    \end{tikzpicture}
    }
    \caption{Level-1 Data Flow Diagram detailing internal compliance processes.}
    \label{fig:dfd_level_1}
\end{figure}

\section{Use Cases}

The use cases define the interaction sequences available to authenticated actors, providing deterministic state definitions for platform operations.

\subsection{Use Case Diagram}

Figure \ref{fig:use_case_diagram} maps the operational workflows orchestrated by the System Administrator and the Security Auditor, establishing boundaries of execution based on the platform's internal RBAC implementation.

\begin{figure}[H]
    \centering
    \resizebox{\textwidth}{!}{
    \begin{tikzpicture}[
        node distance=2cm,
        usecase/.style={ellipse, draw=blue!80!black, thick, fill=blue!10, minimum width=4cm, minimum height=1.2cm, align=center, font=\bfseries\small},
        authcase/.style={ellipse, draw=red!80!black, thick, fill=red!10, minimum width=4cm, minimum height=1.2cm, align=center, font=\bfseries\small},
        stickman/.style={circle, draw=black!80, thick, minimum size=6mm, fill=white,
            append after command={
                \pgfextra{
                    \draw[thick, draw=black!80] (\tikzlastnode.south) -- ++(0,-8mm) coordinate (body_bottom);
                    \draw[thick, draw=black!80] (\tikzlastnode.south) ++(0,-2mm) -- ++(-4mm,-4mm);
                    \draw[thick, draw=black!80] (\tikzlastnode.south) ++(0,-2mm) -- ++(4mm,-4mm);
                    \draw[thick, draw=black!80] (body_bottom) -- ++(-4mm,-6mm);
                    \draw[thick, draw=black!80] (body_bottom) -- ++(4mm,-6mm);
                }
            }
        },
        assoc/.style={thick, draw=black!50},
        include/.style={-Stealth, dashed, thick, draw=black!80},
        inclabel/.style={midway, fill=gray!5, font=\scriptsize, inner sep=2pt, text=black!90}
    ]
        % System Boundary
        \node[draw=gray!80!black, thick, minimum width=12cm, minimum height=13.5cm, fill=gray!5, rounded corners] (system) at (4.5, 0) {};
        \node[anchor=north, font=\bfseries\Large, yshift=-3mm, text=black!80] at (system.north) {ICMS Platform};

        % Actors (Left, well spaced to prevent line clipping)
        \node[stickman] (superadmin) at (-5, 5) {}; 
        \node[align=center, font=\bfseries\small, text=black!90] at (-5, 3.2) {Superadmin};

        \node[stickman] (admin) at (-5, 0) {}; 
        \node[align=center, font=\bfseries\small, text=black!90] at (-5, -1.8) {System\\Administrator};

        \node[stickman] (auditor) at (-5, -5) {}; 
        \node[align=center, font=\bfseries\small, text=black!90] at (-5, -6.8) {Security\\Auditor};

        % Primary Use Cases (Center)
        \node[usecase] (uc1) at (1.5, 5) {Manage Users};
        \node[usecase] (uc2) at (1.5, 3) {Manage Frameworks};
        \node[usecase] (uc3) at (1.5, 1) {Map SIEM Rules};
        \node[usecase] (uc4) at (1.5, -1) {Sync Telemetry};
        \node[usecase] (uc5) at (1.5, -3) {View Dashboard};
        \node[usecase] (uc6) at (1.5, -5) {Export Reports};

        % Included Auth Use Case (Right)
        \node[authcase] (ucauth) at (8, 0) {User Authentication};

        % Actor Associations
        % Superadmin
        \draw[assoc] (superadmin.east) -- (uc1.west);
        \draw[assoc] (superadmin.east) -- (uc2.west);
        \draw[assoc] (superadmin.east) -- (uc3.west);
        \draw[assoc] (superadmin.east) -- (uc4.west);
        \draw[assoc] (superadmin.east) -- (uc5.west);
        \draw[assoc] (superadmin.east) -- (uc6.west);

        % System Admin
        \draw[assoc] (admin.east) -- (uc2.west);
        \draw[assoc] (admin.east) -- (uc3.west);
        \draw[assoc] (admin.east) -- (uc4.west);
        \draw[assoc] (admin.east) -- (uc5.west);
        \draw[assoc] (admin.east) -- (uc6.west);

        % Auditor
        \draw[assoc] (auditor.east) -- (uc5.west);
        \draw[assoc] (auditor.east) -- (uc6.west);

        % Includes (using precise degree anchors to prevent overlap)
        \draw[include] (uc1.east) -- (ucauth.145) node[inclabel, sloped] {$\langle\langle$include$\rangle\rangle$};
        \draw[include] (uc2.east) -- (ucauth.155) node[inclabel, sloped] {$\langle\langle$include$\rangle\rangle$};
        \draw[include] (uc3.east) -- (ucauth.165) node[inclabel, sloped] {$\langle\langle$include$\rangle\rangle$};
        \draw[include] (uc4.east) -- (ucauth.195) node[inclabel, sloped] {$\langle\langle$include$\rangle\rangle$};
        \draw[include] (uc5.east) -- (ucauth.205) node[inclabel, sloped] {$\langle\langle$include$\rangle\rangle$};
        \draw[include] (uc6.east) -- (ucauth.215) node[inclabel, sloped] {$\langle\langle$include$\rangle\rangle$};
        
    \end{tikzpicture}
    }
    \caption{System Use Case Diagram}
    \label{fig:use_case_diagram}
\end{figure}

\subsection{Descriptive Use Cases}

The following tabular specifications define the algorithmic flows, authentication prerequisites, and resulting persistence states for each operation.

\renewcommand{\arraystretch}{1.5}
\begin{table}[H]
    \centering
    \caption{Use Case - Authentication}\label{tab:uc_01}
    \begin{tabular}{|>{\raggedright\arraybackslash}p{3cm}|p{11.5cm}|}
        \hline
        \textbf{Attribute} & \textbf{Description}\\
        \hline
        \textbf{Use Case ID} & UC-01 \\
        \hline
        \textbf{Use Case Name} & User Authentication \\
        \hline
        \textbf{Primary} \par \textbf{Actor} & System Administrator, Security Auditor \\
        \hline
        \textbf{Preconditions} & The user account is registered and active in the system. \\
        \hline
        \textbf{Main Flow} & 1. The user enters their credentials into the login interface.\newline 2. The system verifies the submitted credentials.\newline 3. The system grants access and displays the default dashboard view. \\
        \hline
        \textbf{Alternative Flows} & Invalid credentials \textrightarrow\ The system denies access and displays an authentication error message.\newline Session timeout \textrightarrow\ The system automatically refreshes the session without interrupting the user. \\
                \hline
        \textbf{Frequency of Use} & High (Initiated upon every new user session or token expiry). \\
        \hline
        \textbf{Priority} & Critical (Gatekeeper for all subsequent RBAC routing). \\
        \hline
        \textbf{Exception} & Account locked due to repeated invalid attempts \textrightarrow\ Deny access, log event to security audit trail, and prompt system administrator intervention. \\
\hline
        \textbf{Postconditions} & The user is successfully authenticated and an active session is established. \\
        \hline
    \end{tabular}
\end{table}

\begin{table}[H]
    \centering
    \caption{Use Case - Manage Frameworks}\label{tab:uc_02}
    \begin{tabular}{|>{\raggedright\arraybackslash}p{3cm}|p{11.5cm}|}
        \hline
        \textbf{Attribute} & \textbf{Description}\\
        \hline
        \textbf{Use Case ID} & UC-02 \\
        \hline
        \textbf{Use Case Name} & Manage Wazuh Telemetry Synchronization \\
        \hline
        \textbf{Primary} \par \textbf{Actor} & System Administrator \\
        \hline
        \textbf{Preconditions} & The Administrator is authenticated. $\langle\langle$include$\rangle\rangle$ User Authentication. The external SIEM system is reachable. \\
        \hline
        \textbf{Main Flow} & 1. The Administrator triggers a manual synchronization of security telemetry.\newline 2. The system requests audit data from the external SIEM.\newline 3. The system processes and maps the received data against the linked controls.\newline 4. The system updates the compliance status for each control.\newline 5. The system stores the audit results and notifies the Administrator. \\
        \hline
        \textbf{Alternative Flows} & SIEM system is unreachable \textrightarrow\ The system retries the connection and alerts the Administrator if it ultimately fails. \\
                \hline
        \textbf{Frequency of Use} & High (Continuous via Celery beat background tasks; occasional manual ad-hoc triggers by Superadmin). \\
        \hline
        \textbf{Priority} & Critical (Forms the core data ingestion pipeline for the compliance engine). \\
        \hline
        \textbf{Exception} & Wazuh REST API connection timeout or TLS handshake failure \textrightarrow\ Abort sync, log connection error, and queue retry via Celery exponential backoff. \\
\hline
        \textbf{Postconditions} & The system reflects the latest compliance state based on the synchronized SIEM data. \\
        \hline
    \end{tabular}
\end{table}

\begin{table}[H]
    \centering
    \caption{Use Case - Map Controls}\label{tab:uc_03}
    \begin{tabular}{|>{\raggedright\arraybackslash}p{3cm}|p{11.5cm}|}
        \hline
        \textbf{Attribute} & \textbf{Description}\\
        \hline
        \textbf{Use Case ID} & UC-03 \\
        \hline
        \textbf{Use Case Name} & Map Framework Controls to Wazuh Rules \\
        \hline
        \textbf{Primary} \par \textbf{Actor} & System Administrator \\
        \hline
        \textbf{Preconditions} & The Administrator is authenticated. $\langle\langle$include$\rangle\rangle$ User Authentication. Frameworks and controls exist in the system. \\
        \hline
        \textbf{Main Flow} & 1. The Administrator selects a SIEM rule and links it to a specific compliance control.\newline 2. The system validates the relationship between the rule and the control.\newline 3. The system saves the mapping configuration.\newline 4. The system confirms the successful linkage to the Administrator. \\
        \hline
        \textbf{Alternative Flows} & Invalid control selected \textrightarrow\ The system rejects the mapping and prompts the Administrator to select a valid control. \\
                \hline
        \textbf{Frequency of Use} & Low/Medium (Executed during initial platform setup or when regulatory frameworks release new versions). \\
        \hline
        \textbf{Priority} & High (Defines the relational logic for the entire compliance evaluation). \\
        \hline
        \textbf{Exception} & Mapping collision or invalid/deprecated Wazuh Rule ID \textrightarrow\ Prevent database commit, display strict validation error to the Superadmin. \\
\hline
        \textbf{Postconditions} & A linkage is established between the SIEM rule and the compliance control for future audits. \\
        \hline
    \end{tabular}
\end{table}

\begin{table}[H]
    \centering
    \caption{Use Case - Sync Telemetry}\label{tab:uc_04}
    \begin{tabular}{|>{\raggedright\arraybackslash}p{3cm}|p{11.5cm}|}
        \hline
        \textbf{Attribute} & \textbf{Description}\\
        \hline
        \textbf{Use Case ID} & UC-04 \\
        \hline
        \textbf{Use Case Name} & Monitor Endpoint Compliance Status \\
        \hline
        \textbf{Primary} \par \textbf{Actor} & System Administrator (Trigger) / Celery (System Actor) \\
        \hline
        \textbf{Preconditions} & The user is authenticated. $\langle\langle$include$\rangle\rangle$ User Authentication. \\
        \hline
        \textbf{Main Flow} & 1. The user navigates to the dashboard view.\newline 2. The system aggregates the latest compliance data.\newline 3. The system visualizes the compliance score and detailed control statuses. \\
        \hline
        \textbf{Alternative Flows} & No data available \textrightarrow\ The system displays an empty state with a prompt to synchronize telemetry. \\
                \hline
        \textbf{Frequency of Use} & High (Serves as the primary operational dashboard for continuous monitoring). \\
        \hline
        \textbf{Priority} & High (Core visibility and monitoring function). \\
        \hline
        \textbf{Exception} & Endpoint agent disconnected or offline \textrightarrow\ Mark endpoint status as 'Stale/Offline', visually highlight the disconnection in the UI, and suspend compliance calculations for that node. \\
\hline
        \textbf{Postconditions} & The user is presented with an up-to-date visual compliance posture. \\
        \hline
    \end{tabular}
\end{table}

\begin{table}[H]
    \centering
    \caption{Use Case - View Dashboard}\label{tab:uc_05}
    \begin{tabular}{|>{\raggedright\arraybackslash}p{3cm}|p{11.5cm}|}
        \hline
        \textbf{Attribute} & \textbf{Description}\\
        \hline
        \textbf{Use Case ID} & UC-05 \\
        \hline
        \textbf{Use Case Name} & Resolve Compliance Anomalies \\
        \hline
        \textbf{Primary} \par \textbf{Actor} & Security Auditor, System Administrator \\
        \hline
        \textbf{Preconditions} & The Auditor or Admin is authenticated. Compliance failures have been detected in the system. \\
        \hline
        \textbf{Main Flow} & 1. The user identifies a non-compliant control on the dashboard.\newline 2. The user initiates a resolution workflow and applies corrective actions to the endpoint.\newline 3. The system verifies the change and logs the remediation status. \\
        \hline
        \textbf{Alternative Flows} & Resolution fails \textrightarrow\ The system maintains the 'Failed' state and prompts for manual investigation. \\
                \hline
        \textbf{Frequency of Use} & Medium (Event-driven, based on incoming alerts and failed states). \\
        \hline
        \textbf{Priority} & Medium/High (Critical for operational remediation, but dependent on external IT actions). \\
        \hline
        \textbf{Exception} & Insufficient RBAC permissions to apply resolution (e.g., Auditor attempting Admin action) \textrightarrow\ Deny action, return HTTP 403 Forbidden, and log unauthorized attempt. \\
\hline
        \textbf{Postconditions} & The compliance anomaly is resolved and recorded in the historical audit ledger. \\
        \hline
    \end{tabular}
\end{table}

\begin{table}[H]
    \centering
    \caption{Use Case - Export Audit Reports}\label{tab:uc_06}
    \begin{tabular}{|>{\raggedright\arraybackslash}p{3cm}|p{11.5cm}|}
        \hline
        \textbf{Attribute} & \textbf{Description}\\
        \hline
        \textbf{Use Case ID} & UC-06 \\
        \hline
        \textbf{Use Case Name} & Export Audit Reports \\
        \hline
        \textbf{Primary} \par \textbf{Actor} & Security Auditor \\
        \hline
        \textbf{Preconditions} & The Auditor is authenticated. $\langle\langle$include$\rangle\rangle$ User Authentication. Compliance data exists. \\
        \hline
        \textbf{Main Flow} & 1. The Auditor requests a formal compliance report for a specific framework.\newline 2. The system compiles the historical audit data.\newline 3. The system generates a formatted PDF document.\newline 4. The system prompts the Auditor to download the document. \\
        \hline
        \textbf{Alternative Flows} & Generation failure \textrightarrow\ The system displays a generation error message and logs the failure. \\
                \hline
        \textbf{Frequency of Use} & Low (Triggered periodically based on audit cycles, e.g., monthly or quarterly). \\
        \hline
        \textbf{Priority} & High (Primary executive deliverable for external governance bodies). \\
        \hline
        \textbf{Exception} & PostgreSQL data retrieval timeout or PDF template rendering failure \textrightarrow\ Abort generation, display generic error modal, and log backend stack trace. \\
\hline
        \textbf{Postconditions} & A timestamped audit report is delivered to the Auditor. \\
        \hline

    \end{tabular}
\end{table}
